% Part: incompleteness
% Chapter: incompleteness-provability
% Section: provability-conditions

\documentclass[../../../include/open-logic-section]{subfiles}

\begin{document}

\olfileid{inc}{inp}{prc}

\olsection{شروط \usetoken{S}{derivability} لـ$\Th{PA}$}

حساب بيانو، أو~$\Th{PA}$، هو النظرية التي توسع~$\Th{Q}$ ببديهيات
استقراء لجميع الـ!!{formula}s. وبعبارة أخرى، نضيف إلى~$\Th{Q}$
بديهيات من الصورة
\[
(!A(0) \land \lforall[x][(!A(x) \lif !A(x'))]) \lif \lforall[x][!A(x)]
\]
لكل !!{formula}~$!A$. ولاحظ أن هذا \emph{مخطط} حقًا؛ أي إنه عدد لا
نهائي من البديهيات (وقد تبين أن $\Th{PA}$ {\em ليست} قابلة للمحورة
بمجموعة منتهية). لكن لما كان من الممكن أن نقرر فعليًا ما إذا كانت
سلسلة من الرموز حالة من بديهية استقراء، فإن مجموعة بديهيات~$\Th{PA}$
قابلة للحساب. و$\Th{PA}$ نظرية أمتن بكثير من~$\Th{Q}$. فعلى سبيل
المثال، يمكن بسهولة إثبات إبدالية الجمع والضرب باستعمال الاستقراء على
النحو المألوف. بل يمكن إجراء معظم الحجج التناهية في نظرية الأعداد
والتوافقيات داخل~$\Th{PA}$.

لما كانت $\Th{PA}$ مؤكسمة حسابيًا، فإن محمول !!{derivability}
$\Prf[\Th{PA}](x,y)$ قابل للحساب، ومن ثم ممثل في~$\Th{Q}$ (وبالتالي
في~$\Th{PA}$). وكما سبق، سنجعل $\OPrf[\Th{PA}](x,y)$ يرمز إلى الصيغة
التي تمثل العلاقة. ولتكن $\OProv[\Th{PA}](y)$ الصيغة
$\lexists[x][\OPrf[\Th{PA}](x,y)]$، التي تقول حدسيًا: «إن $y$
!!{derivable} من بديهيات~$\Th{PA}$». وسبب حاجتنا إلى أكثر قليلًا من
بديهيات~$\Th{Q}$ هو أننا نحتاج إلى معرفة أن النظرية التي نستعملها
قوية بما يكفي لكي تتمكن من !!{derive} حقائق أساسية قليلة عن محمول
!!{derivability} هذا. وما نحتاج إليه في الواقع هو الحقائق الآتية:
\begin{enumerate}
\item[P1.] إذا كان $\Th{PA}\Proves !A$، فإن
  $\Th{PA}\Proves\OProv[\Th{PA}](\gn{!A})$.
\item[P2.] لجميع الـ!!{formula}s $!A$ و$!B$،
  \[
  \Th{PA} \Proves \OProv[\Th{PA}](\gn{!A \lif !B}) \lif
  (\OProv[\Th{PA}](\gn{!A}) \lif \OProv[\Th{PA}](\gn{!B})).
  \]
\item[P3.] لكل !!{formula}~$!A$،
  \[
  \Th{PA} \Proves \OProv[\Th{PA}](\gn{!A})
  \lif \OProv[\Th{PA}](\gn{\OProv[\Th{PA}](\gn{!A})}).
  \]
\end{enumerate}
لا سبيل إلى التحقق من تحقق هذه الخواص الثلاث إلا بوصف الـ!!{formula}
$\OProv[\Th{PA}](y)$ بعناية، واستعمال بديهيات~$\Th{PA}$ لوصف
الـ!!{derivation}s الصورية ذات الصلة. الشرطان (1) و(2) سهلان؛ والشرط
(3) هو الذي يتطلب عملًا حقًا. (فكر في نوع العمل الذي يستلزمه ذلك
\dots.) وسيكون تنفيذ التفاصيل مضجرًا وغير ممتع، ولذلك سنطلب منك هنا
أن تقبل أن لـ$\Th{PA}$ الخواص الثلاث المذكورة أعلاه. وسيحقق اختيار
معقول لـ$\OProv[\Th{PA}](y)$ أيضًا
\begin{enumerate}
\item[P4.] إذا كان
  $\Th{PA}\Proves\OProv[\Th{PA}](\gn{!A})$، فإن
  $\Th{PA}\Proves !A$.
\end{enumerate}
لكننا لن نحتاج إلى هذه الحقيقة.

\begin{digress}
وبالمناسبة، كان غودل متكاسلًا بالطريقة نفسها التي نتكاسل بها الآن.
ففي ختام بحثه لعام 1931 يرسم برهان مبرهنة عدم الاكتمال الثانية، ويعد
بتقديم التفاصيل في بحث لاحق. ولم يعد قط إلى ذلك؛ إذ آمن كل من فهم
الحجة بإمكان تنفيذها (فلم يكن بحاجة إلى ملء التفاصيل.)
\end{digress}

\end{document}
